\documentclass[10pt, aspectratio=169]{beamer}
\usepackage{../beamer_theme_408}
\title{408考研零基础 --- 计算机网络}
\subtitle{4-6 数据链路层功能}
\author{}
\date{}
\begin{document}
\frame{\titlepage}
\begin{frame}{本节目标}
    \begin{enumerate}
        \item 掌握数据链路层的五大功能
        \item 理解封装成帧和透明传输
        \item 掌握CRC循环冗余校验原理
        \item 区分检错与纠错
        \item 理解滑动窗口机制的基本概念
    \end{enumerate}
\end{frame}

\begin{frame}{数据链路层功能概述}
    \begin{block}{定位}
        数据链路层位于\textbf{物理层之上、网络层之下}，负责在相邻节点间的链路上可靠地传输数据。
    \end{block}
    \vspace{0.3em}
    \textbf{五大功能}：
    \begin{enumerate}
        \item \textbf{封装成帧}：将网络层的数据报加上首部和尾部，形成帧
        \item \textbf{透明传输}：确保数据中任何比特组合都能正常传输
        \item \textbf{差错控制}：检测和纠正传输过程中的比特错误
        \item \textbf{流量控制}：协调发送方和接收方的速率
        \item \textbf{链路管理}：建立、维护和释放数据链路
    \end{enumerate}
    \vspace{0.3em}
    \begin{exampleblock}{类比}
        数据链路层相当于快递的\textbf{分拣和运输环节}，负责将包裹（帧）从一个站点送到相邻站点。
    \end{exampleblock}
\end{frame}

\begin{frame}{封装成帧}
    \begin{block}{帧结构}
        $$\boxed{\text{帧首部} \;|\; \text{数据部分（网络层PDU）} \;|\; \text{帧尾部}}$$
    \end{block}
    \vspace{0.3em}
    \begin{columns}[T]
        \column{0.5\textwidth}
        \textbf{帧的组成}：
        \begin{itemize}
            \item \textbf{首部}：包含目的地址、源地址等控制信息
            \item \textbf{数据}：上层（网络层）传递下来的数据报
            \item \textbf{尾部}：包含校验码（FCS）等
        \end{itemize}
        \column{0.5\textwidth}
        \textbf{帧定界符}：
        \begin{itemize}
            \item \textbf{SOH}（Start of Header）：帧首定界符
            \item \textbf{EOT}（End of Transmission）：帧尾定界符
            \item 接收方根据定界符识别帧的边界
        \end{itemize}
    \end{columns}
    \vspace{0.3em}
    \begin{block}{最大传输单元（MTU）}
        每种链路层协议规定了帧数据部分的\textbf{最大长度}，超过MTU的数据报需要分片。
    \end{block}
\end{frame}

\begin{frame}{透明传输}
    \begin{block}{问题}
        如果数据部分出现了与\textbf{帧定界符}相同的比特组合，接收方会误认为是帧的边界，导致错误。
    \end{block}
    \vspace{0.3em}
    \begin{columns}[T]
        \column{0.5\textwidth}
        \textbf{解决方法：字节填充法}
        \begin{itemize}
            \item 发送方在数据中每个SOH/EOT前插入\textbf{转义字符（ESC）}
            \item 接收方识别后删除转义字符
            \item 如果数据本身有ESC，则再插入一个ESC
            \item 即：\textbf{ESC + 定界符} 表示数据中的定界符
        \end{itemize}
        \column{0.5\textwidth}
        \begin{exampleblock}{示例}
            发送数据：\texttt{A SOH B}\\
            实际发送：\texttt{A ESC SOH B}\\
            接收方收到 \texttt{ESC SOH} 后，去掉 ESC，恢复为 \texttt{SOH}
        \end{exampleblock}
    \end{columns}
    \vspace{0.3em}
    \begin{alertblock}{透明传输的本质}
        对上层而言，数据链路层提供了"透明"的传输通道，数据中任何比特组合都能通过。
    \end{alertblock}
\end{frame}

\begin{frame}{差错控制 --- 检错与纠错}
    \begin{columns}[T]
        \column{0.5\textwidth}
        \textbf{检错（Detection）}
        \begin{itemize}
            \item 只能检测出是否有错误
            \item 无法定位错误位置
            \item 发现错误后请求重传
            \item 典型方法：\textbf{CRC循环冗余校验}
            \item 原理：用生成多项式计算冗余码
        \end{itemize}
        \column{0.5\textwidth}
        \textbf{纠错（Correction）}
        \begin{itemize}
            \item 不仅检测错误，还能\underline{定位和纠正}错误
            \item 需要更多的冗余信息
            \item 典型方法：\textbf{海明码}
            \item 能纠正一位错误
            \item 开销比检错大
        \end{itemize}
    \end{columns}
    \vspace{0.5em}
    \begin{block}{考研对比}
        \begin{itemize}
            \item CRC：\textbf{检错}能力强，不能纠错，硬件实现简单
            \item 海明码：可以\textbf{纠错}，但要增加更多校验位
        \end{itemize}
    \end{block}
\end{frame}

\begin{frame}{CRC循环冗余校验}
    \begin{block}{基本思想}
        发送方在数据末尾添加\textbf{冗余码（FCS）}，使整个数据能被生成多项式整除；接收方用同一多项式除，余数为0则无错。
    \end{block}
    \vspace{0.3em}
    \textbf{计算步骤}：
    \begin{enumerate}
        \item 确定生成多项式 $P$（如 CRC-16）
        \item 在原数据后添加 $r$ 个0（$r$ = 多项式最高次数）
        \item 用\textbf{模2除法}（异或运算）计算余数
        \item 余数即为 \textbf{FCS}（帧校验序列）
        \item 发送方发送：\textbf{原数据 + FCS}
    \end{enumerate}
    \vspace{0.3em}
    \begin{block}{接收方检错}
        接收方用同一多项式除"数据+FCS"，若余数为0 $\to$ 无错；余数不为0 $\to$ 有错，丢弃或请求重传。
    \end{block}
\end{frame}

\begin{frame}{CRC计算示例}
    \begin{exampleblock}{例题}
        数据 $M = 101001$，生成多项式 $P = 1101$（$r=3$），求FCS？\\
        \textbf{解}：
        \begin{enumerate}
            \item $M$ 后加3个0：$101001000$
            \item 模2除法：
            \begin{align*}
                &\phantom{=}101001000 \div 1101 \\
                &\text{余数} = 001
            \end{align*}
            \item FCS = 001
            \item 发送的码字 = $101001001$
        \end{enumerate}
    \end{exampleblock}
    \vspace{0.3em}
    \begin{exampleblock}{验算}
        接收方用 $101001001 \div 1101$，余数为0 $\to$ 无错误。
    \end{exampleblock}
    \begin{block}{注意}
        模2除法的本质是\textbf{异或}运算，不进位也不借位。
    \end{block}
\end{frame}

\begin{frame}{流量控制 --- 滑动窗口}
    \begin{block}{为什么需要流量控制？}
        发送方的发送速率可能超过接收方的接收能力，导致数据丢失。流量控制用来协调双方速率。
    \end{block}
    \vspace{0.3em}
    \begin{columns}[T]
        \column{0.5\textwidth}
        \textbf{滑动窗口机制}
        \begin{itemize}
            \item \textbf{发送窗口}：发送方可以发送但尚未确认的帧的序号范围
            \item \textbf{接收窗口}：接收方允许接收的帧的序号范围
            \item 窗口大小动态调整
            \item 每收到一个确认，窗口向前滑动
        \end{itemize}
        \column{0.5\textwidth}
        \begin{block}{滑动窗口类型}
            \begin{itemize}
                \item \textbf{停止-等待协议}：发送窗口=$1$，接收窗口=$1$
                \item \textbf{后退N帧协议（GBN）}：发送窗口$>1$，接收窗口=$1$
                \item \textbf{选择重传协议（SR）}：发送窗口$>1$，接收窗口$>1$
            \end{itemize}
        \end{block}
    \end{columns}
    \vspace{0.3em}
    \begin{exampleblock}{核心概念}
        发送窗口大小决定了\textbf{已发送但未确认}的帧的最大数量。
    \end{exampleblock}
\end{frame}

\begin{frame}{本节总结}
    \begin{enumerate}
        \item 数据链路层五大功能：\textbf{封装成帧、透明传输、差错控制、流量控制、链路管理}
        \item \textbf{封装成帧}：帧 = 首部 + 数据 + 尾部，用SOH/EOT定界
        \item \textbf{透明传输}：数据中定界符前加转义字符ESC
        \item \textbf{差错控制}：检错（CRC）检测错误，纠错（海明码）定位并纠正
        \item \textbf{CRC}：模2除法计算FCS，余数为0则无错
        \item \textbf{流量控制}：滑动窗口机制（发送窗口和接收窗口）
    \end{enumerate}
    \vspace{1em}
    \begin{center}
        \textcolor{blue}{\textbf{下节预告：4-7 组帧与差错检测}}
    \end{center}
\end{frame}
\end{document}
